% -------------------------------------------------------------------------------
% Establish page structure & font.
\documentclass[12pt]{report}

\usepackage[total={6.5in, 9in},
	left=1in,
	right=1in,
	top=1in,
	bottom=1in,]{geometry} % Page structure

\usepackage{graphicx} % Required for inserting images
\graphicspath{{../.images/}} % Any additional images I use (BCU logo, etc) are from here.

\usepackage[utf8]{inputenc} % UTF-8 encoding
\usepackage[T1]{fontenc} % T1 font
\usepackage{float}  % Allows for floats to be positioned using [H], which correctly
                    % positions them relative to their location within my LaTeX code.
\usepackage{subcaption}
\usepackage{csquotes}

% -------------------------------------------------------------------------------
% Declare biblatex with custom Harvard BCU styling for referencing.
\usepackage[
    useprefix=true,
    maxcitenames=3,
    maxbibnames=99,
    style=authoryear,
    dashed=false, 
    natbib=true,
    url=false,
    backend=biber
]{biblatex}

\usepackage[british]{babel}

% Additional styling options to ensure Harvard referencing format.
\renewbibmacro*{volume+number+eid}{
    \printfield{volume}
    \setunit*{\addnbspace}
    \printfield{number}
    \setunit{\addcomma\space}
    \printfield{eid}}
\DeclareFieldFormat[article]{number}{\mkbibparens{#1}}

\addbibresource{Report.bib}

% -------------------------------------------------------------------------------
% To prevent "Chapter N" display for each chapter
\usepackage[compact]{titlesec}
\usepackage{wasysym}
\usepackage{import}

\titlespacing*{\chapter}{0pt}{-2cm}{0.5cm}
\titleformat{\chapter}[display]
{\normalfont\bfseries}{}{0pt}{\Huge}

% -------------------------------------------------------------------------------
% Custom macro to make an un-numbered footnote.

\newcommand\blfootnote[1]{
    \begingroup
    \renewcommand\thefootnote{}\footnote{#1}
    \addtocounter{footnote}{-1}
    \endgroup
}

% -------------------------------------------------------------------------------
% Fancy headers; used to show my name, BCU logo and current chapter for the page.
\usepackage{fancyhdr}
\usepackage{calc}
\pagestyle{fancy}

\setlength\headheight{37pt} % Set custom header height to fit the image.

\renewcommand{\chaptermark}[1]{%
    \markboth{#1}{}} % Include chapter name.


% Lewis Higgins - ID 22133848           [BCU LOGO]                [CHAPTER NAME]
\lhead{Lewis Higgins - ID 22133848~~~~~~~~~~~~~~~\includegraphics[width=1.75cm]{BCU}}
\fancyhead[R]{\leftmark}

% ------------------------------------------------------------------------------
% Used to add PDF hyperlinks for figures and the contents page.

\usepackage{hyperref}

\hypersetup{
    colorlinks=true,
    linkcolor=black,
    filecolor=magenta,
    urlcolor=blue,
    citecolor=black,
}

% ------------------------------------------------------------------------------
\usepackage{xcolor} 
\usepackage{colortbl}
\usepackage{longtable}
\usepackage{amssymb}
% ------------------------------------------------------------------------------
\usepackage{tcolorbox}
\newcommand{\para}{\vspace{7pt}\noindent}
% -------------------------------------------------------------------------------

\title{Untitled CMP6228 Report}
\author{Lewis Higgins - Student ID 22133848}
\date{Unknown, 2025}

% -------------------------------------------------------------------------------

\begin{document}


\makeatletter
\begin{titlepage}
    \begin{center}
        \includegraphics[width=0.7\linewidth]{BCU}\\[4ex]
        {\huge \bfseries CMP6228 - Unknown Assignment}\\[2ex]
        {\large \bfseries  \@title}\\[50ex]
        {\@author}\\[2ex]
        {CMP6228 - Deep Neural Networks}\\[2ex]
        {Module Coordinator: Khalid Ismail}\\[10ex]
    \end{center}
\end{titlepage}
\makeatother
\thispagestyle{empty}
\newpage


% Page counter trick so that the contents page doesn't increment it.
\setcounter{page}{0}

\tableofcontents
\thispagestyle{empty}

\chapter*{Introduction}
\addcontentsline{toc}{chapter}{Introduction}
This module consists of two assignments in a similar fashion to CMP6230. 
By week 8, you must have found and written a \textbf{proposal} (20\%) detailing 
the problem to be solved with deep learning.

\para The second assignment is worth 80\%, and is a full report detailing the design, development
and evaluation of a deep learning model to solve the proposed problem. From this, an immediate 
observation is once again the medical field - image classification of people with some kind of disease 
from an X-ray dataset of some sort. Recall previous data sources - Kaggle, IEEE DataPort, UCI ML. Try 
to avoid the same topics as others if possible, though that's perhaps even less likely than it was in other 
modules if image classification is the main objective.

\para \textbf{Assignment 2's deadline is presumably in Week 12, or maybe even later. 
\textit{A draft version is due in Week 11, however.} } 

\para It is unlikely this module can be completed on your laptop, as TensorFlow 
will be used. Colab is proposed as an option, which could perhaps be feasible enough, though 
that may likely depend on the size of the chosen dataset, of which no specific value 
count was given. (Nor feature count, though this is likely irrelevant in an image classification scenario).
Colab ran epochs \~30-40\% quicker at first glance if the runtime was changed to a GPU.
You should fine-tune a good amount of epochs; more epochs just increase training time, but will they actually 
enhance the model that much or will it just plateau? A 3060 Ti processes twice as fast as Colab's free GPU.

\para Because of this assignment being almost identical to CMP6202 (only difference is the model and 
the fact that you'll likely be doing image classification rather than text),
it's likely that elements of it can be adapted to Assignments 1 and 2. It's likely to be a very math-heavy 
module.

\para The Moodle page for this module is not immediately available; as such, you can't really plan ahead 
for what will be done on a week-by-week basis. Because of this uncertainty, it might be best to prioritise 
this module over CMP6207.

\chapter*{Tech stack}
\begin{itemize}
    \item Previously seen libraries:
    \begin{itemize}
        \item Conda 
        \item Numpy
        \item Scipy (sort-of seen before but not properly)
        \item Scikit-learn 
    \end{itemize}
    \item Pillow
    \begin{itemize}
        \item An image processing library.
    \end{itemize}
    \item h5py
    \begin{itemize}
        \item Appears to be used for serialization of models, like Pickle was in CMP6230.
    \end{itemize}
    \item Keras \& TensorFlow
    \begin{itemize}
        \item Keras apparently uses TensorFlow to maximise efficiency with tensor calculations.
        \item (Or the other way around. Keras is a part of TensorFlow now.)
        \item A tensor is a multidimensional array or matrix.
    \end{itemize}
\end{itemize}

\para Keras has two ways of composing models: Sequential and Functional. \textbf{Only sequential will be 
used in the scope of this module.} 
\vspace{8pt}
\subsubsection{Sequential}
Multiple predefined models are stacked in a linear pipeline like a stack or queue:

% The second line of this seems to input a NumPy array, as there's no second axis.
\begin{verbatim}
    model = Sequential()
    model.add(Dense(N_HIDDEN, input_shape=(784,)))
    model.add(Activation('relu'))
    model.add(Dropout(DROPOUT))
    model.add(Dense(N_HIDDEN))
    model.add(Activation('relu'))
    model.add(Dropout(DROPOUT))
    model.add(Dense(nb_classes))
    model.add(Activation('softmax'))
    model.summary()
\end{verbatim}

\subsubsection{Functional}
Uses an API where complex models can be defined such as DAGs (!!!), models with shared layers, 
or multi-output models. Not much info was given on this, likely requiring your own research.
\textbf{Although it's good to know that these exist, they won't be used in this module.}  


\chapter*{Conclusion}
\addcontentsline{toc}{chapter}{Conclusion}

Overall, something was done\dots

% Interesting trick to instead make the chapter number the letter A for this appendix.
\begingroup
\renewcommand\thechapter{A}
\titleformat{\chapter}[display]
{\normalfont\huge\bfseries}{}{20pt}{\Huge}
\setcounter{section}{0} % Set the section counter back to 0 so earlier chapters don't interfere.
\setcounter{figure}{0} % Set the figure counter back to 0 so earlier figures don't interfere.

\chapter*{Appendix A - Windows Subsystem for Linux}
\addcontentsline{toc}{chapter}{Appendix A - Windows Subsystem for Linux}
\markboth{Appendix A}{}

As of TensorFlow 2.10's release in September 2022, Windows is no longer natively supported with GPU access \autocite{tensorflow_install_nodate}.
To circumvent this issue, the official TensorFlow Docker container can be pulled and used on a Windows host \autocite{tensorflow_install_nodate-1}.
Docker aims to mitigate issues of software incompatibility across devices through containers, which act similarly to virtual machines 
in that they emulate another system with the necessary dependencies on a host. In this specific case, Docker uses the Windows Subsystem for 
Linux to virtualize an Ubuntu client that supports TensorFlow's latest releases with GPU support.

\para \textbf{I'm not going to remove this section until I make a final decision on the following:}
\begin{itemize}
    \item My RTX 3060 Ti seems to work at the same speed, if not slower, than a Colab GPU.
    \item My CPU works over 1000x faster than both, but my research shows that this is true for small datasets and the time needed grows exponentially.
    \item It is possible that the GPU speed issues are due to low VRAM allocation, but this cannot actually be resolved without \textbf{DUAL-BOOTING LINUX.}
    \item Though, is that even necessary? Tensorflow 2.10 is a version that supports the 3060Ti. If half of the things online use completely outdated 
    Python versions, why shouldn't I use a slightly outdated TensorFlow version?
    \item All in all, the decision between Colab and local use will depend entirely on the size of the chosen dataset, though it's likely Colab will win 
    out. 
\end{itemize} 

\addcontentsline{toc}{chapter}{Bibliography}
\printbibliography

\end{document}